\documentclass[11pt,a4paper]{article}

% ---------------- Paquetes ----------------
\usepackage[utf8]{inputenc}
\usepackage[T1]{fontenc}
\usepackage[spanish,es-noquoting]{babel}
\usepackage{lmodern}
\usepackage[margin=2.3cm]{geometry}
\usepackage{xcolor}
\usepackage{listings}
\usepackage{booktabs}
\usepackage{enumitem}
\usepackage{tikz}
\usetikzlibrary{arrows.meta,positioning,automata,shapes.geometric,calc}
\usepackage[most]{tcolorbox}
\usepackage{fancyhdr}
\usepackage[hidelinks]{hyperref}

% ---------------- Colores ----------------
\definecolor{fondo}{RGB}{247,247,250}
\definecolor{clave}{RGB}{0,82,155}
\definecolor{coment}{RGB}{100,120,100}
\definecolor{cadena}{RGB}{163,21,21}
\definecolor{acento}{RGB}{0,110,130}

% ---------------- Listings ----------------
\lstset{
  basicstyle=\ttfamily\small,
  backgroundcolor=\color{fondo},
  keywordstyle=\color{clave}\bfseries,
  commentstyle=\color{coment}\itshape,
  stringstyle=\color{cadena},
  numbers=left, numberstyle=\tiny\color{gray}, numbersep=8pt,
  frame=single, rulecolor=\color{gray!40},
  breaklines=true, showstringspaces=false, tabsize=2,
  xleftmargin=1.5em, framexleftmargin=1.2em,
  literate={á}{{\'a}}1 {é}{{\'e}}1 {í}{{\'i}}1 {ó}{{\'o}}1 {ú}{{\'u}}1
           {Á}{{\'A}}1 {É}{{\'E}}1 {Í}{{\'I}}1 {Ó}{{\'O}}1 {Ú}{{\'U}}1
           {ñ}{{\~n}}1 {¡}{{!`}}1 {¿}{{?`}}1 {Ω}{{$\Omega$}}1
}
\lstdefinelanguage{ini}{
  morecomment=[l]{;}, morestring=[b]",
  morekeywords={platform,board,framework,monitor_speed,build_flags,monitor_rts,monitor_dtr}
}

% ---------------- Cajas ----------------
\newtcolorbox{idea}[1][]{colback=acento!6,colframe=acento,fonttitle=\bfseries,title=#1,arc=2pt,boxrule=0.6pt,before upper={\setlength{\parskip}{5pt}}}

% ---------------- Encabezado ----------------
\pagestyle{fancy}
\fancyhf{}
\lhead{\small Alerta ESP32-S3}
\rhead{\small Explicación del código}
\cfoot{\thepage}

\title{\textbf{Alerta ESP32-S3}\\[4pt]\large Explicación paso a paso del código}
\author{Sara}
\date{\today}

\begin{document}
\shorthandoff{"}
\maketitle
\tableofcontents
\newpage

%==========================================================
\section{¿Qué hace el sistema?}
%==========================================================

El sistema es un botón de alerta portátil. La condición que debe cumplir es:

\begin{enumerate}[leftmargin=*]
  \item Al oprimir el botón por primera vez, el sistema pasa de un estado
        \textbf{estable} a un estado de \textbf{alerta} y el LED se enciende,
        indicándole al portador que el botón fue oprimido.
  \item Si fue un accidente, el portador tiene \textbf{30 segundos} para volver a
        oprimir el mismo botón. En ese caso el sistema regresa a \textbf{estable}
        y el LED se apaga.
  \item Si pasan los 30 segundos sin cancelar, se genera una
        \textbf{notificación al computador por WiFi} y el LED queda encendido.
\end{enumerate}

Además, cada vez que la placa arranca (o se le da \textit{reset}) avisa al
computador que quedó conectada al WiFi, para poder comprobar el enlace.

\subsection{Arquitectura}

\begin{center}
\begin{tikzpicture}[
  bloque/.style={draw,rounded corners,fill=acento!8,minimum width=3cm,minimum height=1.3cm,align=center},
  flecha/.style={-{Stealth[length=3mm]},thick}]
  \node[bloque] (btn) {Botón\\(GPIO4)};
  \node[bloque,below=0.6cm of btn] (led) {LED\\(GPIO2)};
  \node[bloque,right=3.2cm of $(btn)!0.5!(led)$,minimum height=2.8cm] (esp) {\textbf{ESP32-S3}\\[2pt]\small\texttt{main.cpp}};
  \node[bloque,right=2.4cm of esp] (pc) {\textbf{Computador}\\[2pt]\small\texttt{servidor\_pc.py}};
  \draw[flecha] (btn.east) -- (btn.east -| esp.west);
  \draw[flecha] (esp.west |- led.east) -- (led.east);
  \draw[flecha,dashed] (esp) -- node[above]{\small WiFi} node[below]{\small HTTP POST} (pc);
\end{tikzpicture}
\end{center}

Son dos programas que corren en lugares distintos:

\begin{itemize}[leftmargin=*]
  \item \textbf{Firmware} (\texttt{src/main.cpp}): corre \emph{dentro} de la ESP32-S3.
        Lee el botón, controla el LED, lleva el tiempo y envía los mensajes por WiFi.
  \item \textbf{Servidor} (\texttt{pc/servidor\_pc.py}): corre en el computador.
        Espera los mensajes de la ESP32 y los muestra en la terminal.
\end{itemize}

%==========================================================
\section{Hardware y conexiones}
%==========================================================

\begin{center}
\begin{tabular}{@{}llll@{}}
\toprule
\textbf{Componente} & \textbf{Pata} & \textbf{Va a} & \textbf{Nota}\\
\midrule
Pulsador & pata 1 & GPIO4 & \textit{pull-up} interno\\
         & pata 2 & GND   & \\
Resistencia 220\,$\Omega$ & un extremo & GPIO2 & limita la corriente del LED\\
         & otro extremo & ánodo LED (pata larga) & \\
LED      & cátodo (pata corta) & GND & \\
\bottomrule
\end{tabular}
\end{center}

\begin{center}
\begin{tikzpicture}[thick,scale=1]
  % Placa
  \draw[fill=gray!10,rounded corners] (0,0) rectangle (3,4);
  \node at (1.5,3.6) {\textbf{ESP32-S3}};
  \node[anchor=east] at (3,2.8) {\small GPIO4};
  \node[anchor=east] at (3,1.8) {\small GPIO2};
  \node[anchor=east] at (3,0.6) {\small GND};
  \fill (3,2.8) circle (2pt) (3,1.8) circle (2pt) (3,0.6) circle (2pt);
  % Botón (fila superior, baja a GND por la derecha)
  \draw (3,2.8) -- (5.5,2.8);
  \draw (5.5,2.8) circle (1.5pt) (6.5,2.8) circle (1.5pt);
  \draw (5.55,2.85) -- (6.4,3.2);
  \node[above] at (6,3.2) {\small Pulsador};
  \draw (6.55,2.8) -- (9,2.8) -- (9,0.6);
  % Resistencia y LED (fila del medio)
  \draw (3,1.8) -- (3.6,1.8);
  \draw (3.6,1.8) -- (3.7,1.95) -- (3.85,1.65) -- (4,1.95) -- (4.15,1.65) -- (4.3,1.95) -- (4.4,1.8);
  \node[above] at (4,1.95) {\small 220\,$\Omega$};
  \draw (4.4,1.8) -- (6.6,1.8);
  \draw[fill=acento!30] (6.6,1.55) -- (6.6,2.05) -- (7.0,1.8) -- cycle;
  \draw (7.0,1.55) -- (7.0,2.05);
  \draw[-{Stealth[length=1.5mm]},thin] (6.8,2.1) -- (7.05,2.4);
  \draw[-{Stealth[length=1.5mm]},thin] (6.95,2.05) -- (7.2,2.35);
  \node[below] at (6.8,1.55) {\small LED};
  \draw (7.0,1.8) -- (7.6,1.8) -- (7.6,0.6);
  \fill (7.6,0.6) circle (2pt);
  % GND común
  \draw (3,0.6) -- (9,0.6);
\end{tikzpicture}
\end{center}

\begin{idea}[¿Por qué el botón no lleva resistencia?]
El código configura el pin con \texttt{INPUT\_PULLUP}: la ESP32 conecta
internamente el GPIO4 a 3.3\,V a través de una resistencia. Sin presionar, el pin
lee \texttt{HIGH}; al presionar, el botón lo une a GND y lee \texttt{LOW}. Por eso
en el código ``presionado'' significa \texttt{LOW}.
\end{idea}

%==========================================================
\section{Estructura del repositorio}
%==========================================================

\begin{lstlisting}[numbers=none,language={}]
alerta-esp32s3/
|-- platformio.ini          configuración de la placa y del proyecto
|-- src/main.cpp            firmware (lo que corre en la ESP32)
|-- include/
|   |-- config.h            pines y tiempos
|   |-- secrets.example.h   plantilla de WiFi / URL del PC
|   `-- secrets.h           tus datos reales (NO se sube a git)
|-- pc/servidor_pc.py       servidor que corre en el computador
|-- docs/explicacion.tex    este documento
|-- .gitignore              archivos que git ignora
`-- README.md               instrucciones rápidas
\end{lstlisting}

La idea de separar en archivos es que cada cosa tenga su lugar: si cambias un pin
solo tocas \texttt{config.h}; si cambias de red WiFi solo tocas \texttt{secrets.h};
la lógica en \texttt{main.cpp} no se toca.

%==========================================================
\section{\texttt{platformio.ini}: configuración del proyecto}
%==========================================================

PlatformIO lee este archivo para saber para qué placa compilar y cómo.

\begin{lstlisting}[language=ini]
[env:esp32-s3-devkitc-1]
platform = espressif32
board = esp32-s3-devkitc-1
framework = arduino
monitor_speed = 115200
; Para que Serial salga por el puerto USB nativo del S3 (quitar si usas el puerto UART)
build_flags =
    -D ARDUINO_USB_CDC_ON_BOOT=1
; Evita que el monitor reinicie o bloquee la placa por el USB nativo
monitor_rts = 0
monitor_dtr = 0
\end{lstlisting}

\begin{description}[leftmargin=*,font=\ttfamily]
  \item[{[env:...]}] Nombre del ``entorno''. Un proyecto podría tener varios (por
        ejemplo, para dos placas distintas).
  \item[platform] Familia de chips: \texttt{espressif32} son las ESP32.
  \item[board] Modelo exacto de la placa. Define memoria, pines, etc.
  \item[framework] Usamos el entorno de \textbf{Arduino}, por eso existen
        \texttt{setup()}, \texttt{loop()}, \texttt{digitalRead()}, etc.
  \item[monitor\_speed] Velocidad del monitor serial; debe coincidir con
        \texttt{Serial.begin(115200)}.
  \item[build\_flags] \texttt{ARDUINO\_USB\_CDC\_ON\_BOOT=1} hace que los mensajes de
        \texttt{Serial} salgan por el USB nativo del S3 (el COM que aparece como
        ``Dispositivo serie USB'').
  \item[monitor\_rts / monitor\_dtr] Evitan que al abrir el monitor la placa se
        reinicie o quede bloqueada.
\end{description}

%==========================================================
\section{\texttt{config.h} y \texttt{secrets.h}: los datos}
%==========================================================

\subsection{\texttt{config.h}}

\begin{lstlisting}[language=C++]
#pragma once

// ---------- Pines ----------
constexpr int PIN_BOTON = 4;   // botón entre GPIO4 y GND (pull-up interno)
constexpr int PIN_LED   = 2;   // LED + resistencia (220-330 Ω) a GND

// ---------- Tiempos ----------
constexpr unsigned long VENTANA_MS   = 30000;  // 30 s para cancelar la alerta
constexpr unsigned long DEBOUNCE_MS  = 50;
constexpr unsigned long REINTENTO_MS = 3000;   // reintento si falla el envío
\end{lstlisting}

\begin{itemize}[leftmargin=*]
  \item \texttt{\#pragma once}: evita que el archivo se incluya dos veces al compilar.
  \item \texttt{constexpr}: constantes fijas que el compilador conoce desde el inicio.
  \item Los tiempos están en \textbf{milisegundos}: \texttt{30000} ms = 30 s.
  \item \texttt{DEBOUNCE\_MS}: tiempo para filtrar los rebotes del botón (sección \ref{sec:debounce}).
  \item \texttt{REINTENTO\_MS}: cada cuánto se reintenta un envío que falló.
\end{itemize}

\subsection{\texttt{secrets.h}}

\begin{lstlisting}[language=C++]
#pragma once
#define WIFI_SSID "TU_RED"
#define WIFI_PASS "TU_CLAVE"
#define PC_URL    "http://192.168.1.100:5000/alerta"  // IP que muestra servidor_pc.py
\end{lstlisting}

\begin{itemize}[leftmargin=*]
  \item \texttt{WIFI\_SSID} y \texttt{WIFI\_PASS}: nombre y clave de la red. La ESP32
        solo se conecta a redes de \textbf{2.4\,GHz}.
  \item \texttt{PC\_URL}: dirección del servidor. Tiene tres partes: la IP del
        computador, el puerto (\texttt{5000}) y la ruta (\texttt{/alerta}).
  \item Este archivo está en \texttt{.gitignore}, así la clave nunca se sube al repositorio.
        En el repo solo va \texttt{secrets.example.h} como plantilla.
\end{itemize}

%==========================================================
\section{\texttt{main.cpp}: el firmware paso a paso}
%==========================================================

\subsection{Máquina de estados}

Todo el programa se organiza como una \textbf{máquina de estados}: en cada momento
el sistema está en uno de tres estados, y ciertos eventos lo hacen pasar a otro.

\begin{center}
\begin{tikzpicture}[->,>={Stealth[length=3mm]},node distance=4.2cm,thick,
  every state/.style={fill=acento!10,draw=acento,minimum size=2.2cm,align=center}]
  \node[state,initial,initial text={arranque}] (E) {ESTABLE\\\small LED off};
  \node[state,right of=E] (P) {PENDIENTE\\\small LED on};
  \node[state,right of=P] (S) {ENVIADA\\\small LED on};
  \path (E) edge[bend left=25] node[above]{\small botón} (P)
        (P) edge[bend left=25] node[below]{\small botón ($<30$ s)} (E)
        (P) edge node[above]{\small pasan 30 s} (S);
\end{tikzpicture}
\end{center}

\begin{center}
\begin{tabular}{@{}lll@{}}
\toprule
\textbf{Estado} & \textbf{LED} & \textbf{Qué hace}\\
\midrule
\texttt{ESTABLE}   & apagado   & Espera que se oprima el botón.\\
\texttt{PENDIENTE} & encendido & Cuenta 30 s; si se oprime otra vez, cancela.\\
\texttt{ENVIADA}   & encendido & Envía la alerta al PC (reintenta si falla).\\
\bottomrule
\end{tabular}
\end{center}

\subsection{Librerías y archivos incluidos}

\begin{lstlisting}[language=C++]
#include <Arduino.h>
#include <WiFi.h>
#include <HTTPClient.h>

#include "config.h"
#include "secrets.h"
\end{lstlisting}

\begin{itemize}[leftmargin=*]
  \item \texttt{Arduino.h}: funciones básicas (\texttt{pinMode}, \texttt{digitalRead},
        \texttt{millis}, \texttt{Serial}\ldots). En el Arduino IDE se incluye solo; en
        PlatformIO hay que ponerlo.
  \item \texttt{WiFi.h}: conectarse a la red WiFi.
  \item \texttt{HTTPClient.h}: enviar peticiones HTTP (como un navegador) al servidor.
  \item Las comillas \texttt{"..."} indican archivos propios del proyecto (carpeta \texttt{include/}).
\end{itemize}

\subsection{Variables globales}

\begin{lstlisting}[language=C++]
enum Estado { ESTABLE, PENDIENTE, ENVIADA };
Estado estado = ESTABLE;

unsigned long tInicioAlerta = 0;
unsigned long tUltimoIntento = 0;
bool notificado = false;
bool primerIntento = true;
bool avisoConexion = false;      // ya se avisó al PC que la placa está en línea
unsigned long tAviso = 0;

// Debounce
int lecturaAnterior = HIGH;
int estadoBoton = HIGH;
unsigned long tCambio = 0;
\end{lstlisting}

Las variables \textbf{globales} se declaran fuera de cualquier función. Eso
significa que \texttt{setup()}, \texttt{loop()} y todas las demás funciones las
pueden leer y cambiar, y que \textbf{conservan su valor} entre una vuelta y otra
del \texttt{loop()}. Son la ``memoria'' del programa: gracias a ellas la placa
recuerda en qué estado está y cuándo pasó cada cosa.

\begin{idea}[Antes de empezar: ¿qué es \texttt{millis()}?]
\texttt{millis()} devuelve cuántos milisegundos lleva encendida la placa: justo al
arrancar vale 0, al segundo vale 1000, a los 10 segundos vale 10\,000, y así
sigue subiendo. Es como un reloj que nunca se detiene.

Varias variables de abajo guardan un valor de \texttt{millis()} para ``anotar la
hora'' en que pasó algo. Luego, restando \texttt{millis() - hora\_anotada}, se
sabe cuánto tiempo ha pasado desde ese momento.

¿Por qué no usar simplemente \texttt{delay(30000)}? Porque \texttt{delay} congela
la placa: durante esos 30 s no podría leer el botón y nadie podría cancelar la
alerta. Con \texttt{millis()} el programa nunca se detiene.
\end{idea}

\subsubsection*{1. \texttt{enum Estado} y \texttt{estado}: en qué parte del proceso vamos}

\begin{lstlisting}[language=C++,numbers=none]
enum Estado { ESTABLE, PENDIENTE, ENVIADA };
Estado estado = ESTABLE;
\end{lstlisting}

\begin{itemize}[leftmargin=*]
  \item \texttt{enum Estado \{ ... \}} \textbf{no guarda nada todavía}: solo inventa
        un tipo de dato nuevo llamado \texttt{Estado}, que únicamente puede tomar
        tres valores: \texttt{ESTABLE}, \texttt{PENDIENTE} o \texttt{ENVIADA}.
  \item Por dentro, el compilador los convierte en números
        (\texttt{ESTABLE} = 0, \texttt{PENDIENTE} = 1, \texttt{ENVIADA} = 2). Podríamos
        haber usado esos números directamente, pero compare:
        \begin{center}
        \texttt{if (estado == 1)} \quad frente a \quad \texttt{if (estado == PENDIENTE)}
        \end{center}
        El segundo se entiende solo; el primero obliga a recordar qué significaba
        el 1. Además, el compilador no deja asignar un valor que no sea uno de los
        tres, lo que evita errores.
  \item \texttt{Estado estado = ESTABLE;} ahora sí \textbf{crea la variable}
        \texttt{estado} (de tipo \texttt{Estado}) y le da el valor inicial
        \texttt{ESTABLE}. Por eso, cada vez que la placa enciende o se reinicia,
        arranca en reposo con el LED apagado.
  \item A lo largo del programa esta variable se va cambiando
        (\texttt{estado = PENDIENTE;}, \texttt{estado = ENVIADA;}\ldots) y el
        \texttt{switch} del \texttt{loop()} la consulta para decidir qué hacer.
\end{itemize}

\subsubsection*{2. \texttt{tInicioAlerta}: el cronómetro de los 30 segundos}

\begin{lstlisting}[language=C++,numbers=none]
unsigned long tInicioAlerta = 0;
\end{lstlisting}

\begin{itemize}[leftmargin=*]
  \item El nombre significa \emph{tiempo de inicio de la alerta}. Guarda el valor
        de \texttt{millis()} en el instante exacto en que se oprimió el botón.
  \item Es \texttt{unsigned long} porque \texttt{millis()} devuelve números muy
        grandes y nunca negativos (este tipo aguanta hasta unos 49 días).
  \item Ejemplo paso a paso:
        \begin{enumerate}
          \item La placa lleva 12 segundos encendida y se oprime el botón. En ese
                momento \texttt{millis()} vale 12\,000, así que se anota
                \texttt{tInicioAlerta = 12000}.
          \item En cada vuelta del \texttt{loop()} se calcula
                \texttt{millis() - tInicioAlerta}:
                \begin{itemize}
                  \item a los 20\,s de encendida: $20000 - 12000 = 8000$ ms
                        $\rightarrow$ han pasado 8\,s, todavía se puede cancelar;
                  \item a los 42\,s de encendida: $42000 - 12000 = 30000$ ms
                        $\rightarrow$ ya pasaron los 30\,s, se envía la alerta.
                \end{itemize}
        \end{enumerate}
  \item Arranca en 0 porque todavía no se ha oprimido nada; su valor solo importa
        mientras el estado es \texttt{PENDIENTE}.
\end{itemize}

\subsubsection*{3. \texttt{notificado}, \texttt{primerIntento}, \texttt{tUltimoIntento}: enviar la alerta sin rendirse}

\begin{lstlisting}[language=C++,numbers=none]
unsigned long tUltimoIntento = 0;
bool notificado = false;
bool primerIntento = true;
\end{lstlisting}

El envío por WiFi puede fallar (el PC no respondió, la red estaba ocupada\ldots).
La alerta es importante, así que el programa no se rinde: reintenta hasta que el PC
confirme. Para eso necesita recordar tres cosas:

\begin{center}
\begin{tabular}{@{}lll@{}}
\toprule
\textbf{Variable} & \textbf{Tipo} & \textbf{Pregunta que responde}\\
\midrule
\texttt{notificado}     & sí/no (\texttt{bool}) & ¿El PC ya confirmó que recibió la alerta?\\
\texttt{primerIntento}  & sí/no (\texttt{bool}) & ¿Todavía no se ha intentado enviar ni una vez?\\
\texttt{tUltimoIntento} & tiempo (ms)           & ¿A qué hora fue el último intento?\\
\bottomrule
\end{tabular}
\end{center}

Así trabajan juntas (ejemplo):
\begin{enumerate}[leftmargin=*]
  \item Se cumplen los 30\,s. El programa pone \texttt{notificado = false} (aún no
        llega) y \texttt{primerIntento = true} (aún no se ha intentado).
  \item Como \texttt{primerIntento} es \texttt{true}, envía \textbf{de inmediato}.
        Anota la hora en \texttt{tUltimoIntento} (digamos 42\,000) y pone
        \texttt{primerIntento = false}.
  \item \textbf{Si el PC respondió OK}: \texttt{notificado = true}. Listo, no se
        vuelve a enviar nunca más.
  \item \textbf{Si falló}: \texttt{notificado} sigue en \texttt{false}. En las
        siguientes vueltas el programa calcula \texttt{millis() - tUltimoIntento}:
        a los 43\,000 da 1000 (esperar), a los 45\,000 da 3000 $\rightarrow$ ya
        pasaron 3\,s, reintenta y anota la nueva hora.
  \item Se repite el paso 4 cada 3\,s hasta que el PC confirme.
\end{enumerate}

\textbf{¿Para qué sirve \texttt{primerIntento}?} Para que el primer envío salga al
instante en vez de esperar 3\,s. \textbf{¿Y por qué esperar 3\,s entre intentos?}
Porque sin esa pausa la placa intentaría enviar miles de veces por segundo y
saturaría la red.

\subsubsection*{4. \texttt{avisoConexion}, \texttt{tAviso}: avisar que la placa está en línea}

\begin{lstlisting}[language=C++,numbers=none]
bool avisoConexion = false;
unsigned long tAviso = 0;
\end{lstlisting}

Al encender, la placa le manda al PC el mensaje ``ESP32 conectada'' para
confirmar que el WiFi funciona. Este envío también puede fallar, así que usa la
misma idea de reintentos que la alerta, pero con solo dos variables:

\begin{itemize}[leftmargin=*]
  \item \texttt{avisoConexion}: ¿el PC ya confirmó el aviso? Empieza en
        \texttt{false}; cuando el PC responde OK pasa a \texttt{true} y el aviso
        no se vuelve a mandar.
  \item \texttt{tAviso}: a qué hora fue el último intento del aviso. Se usa para
        reintentar cada 3\,s (\texttt{millis() - tAviso >= 3000}).
\end{itemize}

Aquí no hace falta un ``primer intento inmediato'': como \texttt{tAviso} empieza en
0 y la placa tarda varios segundos en arrancar, la resta ya da más de 3000 la
primera vez y el aviso sale enseguida.

Como se reinician con cada encendido, después de cada \textit{reset} el aviso se
vuelve a enviar. Por eso sirve para comprobar que la placa quedó enlazada.

\subsubsection*{5. \texttt{lecturaAnterior}, \texttt{estadoBoton}, \texttt{tCambio}: leer bien el botón}

\begin{lstlisting}[language=C++,numbers=none]
int lecturaAnterior = HIGH;
int estadoBoton = HIGH;
unsigned long tCambio = 0;
\end{lstlisting}

Estas tres son para el \textit{debounce} (filtro de rebotes), que se explica en
detalle en la sección \ref{sec:debounce}. En resumen:

\begin{center}
\begin{tabular}{@{}lp{9.3cm}@{}}
\toprule
\textbf{Variable} & \textbf{Qué guarda}\\
\midrule
\texttt{lecturaAnterior} & Lo que leyó el pin en la vuelta anterior del \texttt{loop()}, aunque sea un rebote.\\
\texttt{estadoBoton}     & El estado del botón ya \textbf{confirmado} (estable por más de 50\,ms).\\
\texttt{tCambio}         & La hora del último cambio en la lectura, para medir si ya lleva 50\,ms quieta.\\
\bottomrule
\end{tabular}
\end{center}

Todas arrancan en \texttt{HIGH} porque, con el \textit{pull-up}, \texttt{HIGH}
significa ``botón suelto'', que es como está el botón al encender.

\subsection{Lectura del botón con \textit{debounce}}\label{sec:debounce}

\begin{lstlisting}[language=C++]
bool botonPresionado() {
  int lectura = digitalRead(PIN_BOTON);
  if (lectura != lecturaAnterior) tCambio = millis();
  lecturaAnterior = lectura;

  if (millis() - tCambio > DEBOUNCE_MS && lectura != estadoBoton) {
    estadoBoton = lectura;
    if (estadoBoton == LOW) return true;  // flanco de bajada = presionado
  }
  return false;
}
\end{lstlisting}

Un botón mecánico no pasa limpio de abierto a cerrado: durante unos milisegundos
``rebota'' y el pin lee \texttt{HIGH, LOW, HIGH, LOW\ldots}. Sin filtro, una sola
presión se contaría como varias y el sistema se activaría y cancelaría solo.

\begin{enumerate}[leftmargin=*]
  \item Se lee el pin (línea 2).
  \item Si la lectura cambió respecto a la anterior, se reinicia el cronómetro
        \texttt{tCambio} (línea 3). Mientras haya rebotes, el cronómetro se sigue
        reiniciando.
  \item Solo cuando la lectura lleva \textbf{más de 50 ms estable} y es distinta al
        último estado aceptado, se acepta el cambio (líneas 6--7).
  \item Si el nuevo estado es \texttt{LOW} (presionado), la función devuelve
        \texttt{true}. Así se detecta el \textbf{momento} de oprimir (flanco de
        bajada), no el tiempo que se mantiene oprimido.
\end{enumerate}

\subsection{Conexión al WiFi}

\begin{lstlisting}[language=C++]
void conectarWiFi() {
  if (WiFi.status() == WL_CONNECTED) return;
  WiFi.begin(WIFI_SSID, WIFI_PASS);
  Serial.print("Conectando WiFi");
  unsigned long t0 = millis();
  while (WiFi.status() != WL_CONNECTED && millis() - t0 < 15000) {
    delay(250);
    Serial.print(".");
  }
  Serial.println(WiFi.status() == WL_CONNECTED ? " OK" : " FALLÓ");
  if (WiFi.status() == WL_CONNECTED) {
    Serial.print("IP: ");
    Serial.println(WiFi.localIP());
  }
}
\end{lstlisting}

Esta función se encarga de que la placa quede conectada a la red. Se llama antes
de cada envío, por si el WiFi se cayó en algún momento.

\begin{description}[leftmargin=*,style=nextline]
  \item[Línea 2 --- ¿ya estoy conectada?]
        \texttt{WiFi.status()} pregunta el estado de la conexión y
        \texttt{WL\_CONNECTED} es la respuesta que significa ``conectada''. Si ya lo
        está, \texttt{return} sale de la función de inmediato: no tiene sentido
        volver a conectar.
  \item[Línea 3 --- empezar a conectar]
        \texttt{WiFi.begin(red, clave)} arranca la conexión con los datos de
        \texttt{secrets.h}. Ojo: esta línea \emph{no espera} a que termine; solo da
        la orden. Conectarse toma unos segundos.
  \item[Líneas 4--5 --- anotar la hora]
        Escribe ``Conectando WiFi'' en el monitor y guarda en \texttt{t0} la hora de
        inicio, para poder medir cuánto lleva intentando.
  \item[Líneas 6--9 --- esperar, pero no para siempre]
        El \texttt{while} se repite \textbf{mientras se cumplan las dos condiciones}
        (\texttt{\&\&} significa ``y''):
        \begin{itemize}
          \item \texttt{WiFi.status() != WL\_CONNECTED}: todavía \textbf{no} está conectada, y
          \item \texttt{millis() - t0 < 15000}: han pasado \textbf{menos} de 15\,s.
        \end{itemize}
        En cada vuelta espera 250\,ms y escribe un punto. En el monitor se ve así:
        \texttt{Conectando WiFi.......}. El ciclo termina en cuanto se conecta
        \textbf{o} se cumplen los 15\,s, lo que pase primero. Sin ese límite, si la
        red no existe, la placa se quedaría aquí atrapada para siempre.
  \item[Línea 10 --- decir cómo le fue]
        Esta línea usa el \textbf{operador ternario}, una forma corta de escribir
        un \texttt{if/else}. Su forma general es:
        \begin{center}
        \texttt{condición ? valor\_si\_es\_verdad : valor\_si\_es\_falso}
        \end{center}
        Entonces \texttt{WiFi.status() == WL\_CONNECTED ? " OK" : " FALLÓ"} vale
        \texttt{" OK"} si se conectó y \texttt{" FALLÓ"} si no. Ese texto se le
        pasa a \texttt{Serial.println}. Es exactamente lo mismo que escribir:
\begin{lstlisting}[language=C++,numbers=none]
if (WiFi.status() == WL_CONNECTED) {
  Serial.println(" OK");
} else {
  Serial.println(" FALLÓ");
}
\end{lstlisting}
        pero en una sola línea. El resultado en el monitor queda
        \texttt{Conectando WiFi....... OK}.
  \item[Líneas 11--14 --- mostrar la IP]
        Solo si se conectó, imprime la dirección IP que el router le asignó a la
        placa. \texttt{Serial.print} escribe sin saltar de línea y
        \texttt{Serial.println} escribe y salta de línea, por eso ambas partes
        quedan juntas: \texttt{IP: 192.168.5.108}. Esa IP identifica a la placa
        dentro de la red, y es la misma que luego muestra el servidor.
\end{description}

\subsection{Conexión con indicador en el LED}

\begin{lstlisting}[language=C++]
void conectarWiFiConLed() {
  WiFi.begin(WIFI_SSID, WIFI_PASS);
  Serial.print("Conectando WiFi");
  unsigned long t0 = millis();
  while (WiFi.status() != WL_CONNECTED && millis() - t0 < 15000) {
    digitalWrite(PIN_LED, !digitalRead(PIN_LED));
    delay(250);
    Serial.print(".");
  }
  if (WiFi.status() == WL_CONNECTED) {
    Serial.printf(" OK\nIP: %s\n", WiFi.localIP().toString().c_str());
    digitalWrite(PIN_LED, HIGH);
    delay(2000);
  } else {
    Serial.println(" FALLÓ");
    for (int i = 0; i < 20; i++) {
      digitalWrite(PIN_LED, !digitalRead(PIN_LED));
      delay(80);
    }
  }
  digitalWrite(PIN_LED, LOW);
}
\end{lstlisting}

Hace lo mismo que \texttt{conectarWiFi()}, pero además usa el LED para mostrar qué
está pasando. Así se puede saber si la placa se conectó \textbf{sin tener el
monitor abierto}, por ejemplo cuando está conectada a un cargador. Solo se usa una
vez, al encender (desde \texttt{setup()}).

\begin{description}[leftmargin=*,style=nextline]
  \item[Líneas 2--4 --- igual que antes]
        Da la orden de conectar, escribe el mensaje y anota la hora de inicio.
  \item[Líneas 5--9 --- parpadeo mientras conecta]
        Es el mismo \texttt{while} de 15\,s como máximo, con una línea nueva:
\begin{lstlisting}[language=C++,numbers=none]
digitalWrite(PIN_LED, !digitalRead(PIN_LED));
\end{lstlisting}
        Se lee de adentro hacia afuera:
        \begin{enumerate}
          \item \texttt{digitalRead(PIN\_LED)}: ¿cómo está el LED ahora? (\texttt{HIGH} = prendido, \texttt{LOW} = apagado).
          \item \texttt{!}: significa ``lo contrario''. \texttt{!HIGH} es \texttt{LOW} y \texttt{!LOW} es \texttt{HIGH}.
          \item \texttt{digitalWrite(PIN\_LED, ...)}: pone el LED en ese valor contrario.
        \end{enumerate}
        Resultado: si estaba prendido lo apaga y si estaba apagado lo prende. Como
        esto se repite cada 250\,ms, el LED \textbf{parpadea} mientras intenta
        conectar.
  \item[Líneas 10--13 --- si se conectó: LED fijo 2 segundos]
        \texttt{Serial.printf} imprime texto con ``huecos'' que se rellenan:
        \begin{itemize}
          \item \texttt{\%s} es un hueco para un texto; se rellena con la IP.
          \item \texttt{\textbackslash n} es un salto de línea.
          \item \texttt{WiFi.localIP().toString().c\_str()} convierte la IP al tipo
                de texto que entiende \texttt{printf}: \texttt{localIP()} da la IP,
                \texttt{toString()} la vuelve texto (\texttt{String}) y
                \texttt{c\_str()} la pasa al formato de texto básico de C.
        \end{itemize}
        En el monitor queda \texttt{ OK} y en la siguiente línea
        \texttt{IP: 192.168.5.108}. Después prende el LED y lo deja así 2\,s
        (\texttt{delay(2000)}): esa es la señal de ``conectada''.
  \item[Líneas 14--20 --- si falló: parpadeo rápido]
        Escribe \texttt{FALLÓ} y usa un ciclo \texttt{for} que se repite 20 veces
        (\texttt{i} va de 0 a 19). Cada vez invierte el LED y espera solo 80\,ms.
        20 cambios equivalen a 10 parpadeos muy rápidos, claramente distintos del
        parpadeo lento de ``conectando''.
  \item[Línea 21 --- apagar]
        Pase lo que pase, al final apaga el LED para empezar en estado
        \texttt{ESTABLE} (LED apagado).
\end{description}

\begin{center}
\begin{tabular}{@{}ll@{}}
\toprule
\textbf{Lo que ves en el LED} & \textbf{Significado}\\
\midrule
Parpadeo lento (cada 250\,ms)   & Conectando al WiFi\\
Fijo 2 segundos y se apaga      & Conectada\\
Parpadeo muy rápido y se apaga  & No se pudo conectar (revisar red o clave)\\
\bottomrule
\end{tabular}
\end{center}

\subsection{Envío de eventos al computador}

\begin{lstlisting}[language=C++]
bool enviarEvento(const char* evento) {
  conectarWiFi();
  if (WiFi.status() != WL_CONNECTED) return false;

  HTTPClient http;
  http.setConnectTimeout(2000);
  http.setTimeout(2000);
  http.begin(PC_URL);
  http.addHeader("Content-Type", "application/json");
  String body = String("{\"dispositivo\":\"ESP32-S3\",\"evento\":\"") + evento +
                "\",\"ip\":\"" + WiFi.localIP().toString() +
                "\",\"mac\":\"" + WiFi.macAddress() + "\"}";
  int code = http.POST(body);
  http.end();

  Serial.printf("Evento %s enviado, respuesta HTTP: %d\n", evento, code);
  return code == 200;
}

bool enviarNotificacion() { return enviarEvento("ALERTA"); }
\end{lstlisting}

Esta función es la que ``habla'' con el computador. Recibe el nombre del evento
(\texttt{"CONECTADA"} o \texttt{"ALERTA"}), lo envía y devuelve
\texttt{true} si el PC confirmó que lo recibió o \texttt{false} si no.

\begin{idea}[¿Qué es HTTP y qué es un POST?]
HTTP es el mismo idioma que usa un navegador para pedir páginas web. Hay varios
tipos de mensajes; un \textbf{POST} significa ``te envío estos datos''. La ESP32
actúa como un navegador que envía un formulario, y el servidor de Python actúa
como la página web que lo recibe y contesta. La respuesta trae un
\textbf{código}: \texttt{200} significa ``recibido OK'', \texttt{404} ``esa
dirección no existe'', y un número negativo significa que ni siquiera se pudo
conectar.
\end{idea}

\begin{description}[leftmargin=*,style=nextline]
  \item[Línea 1 --- la firma de la función]
        \texttt{bool} indica que devuelve verdadero o falso.
        \texttt{const char* evento} es el texto que se le pasa (por ejemplo
        \texttt{"ALERTA"}); \texttt{const} indica que la función no lo modifica.
  \item[Líneas 2--3 --- asegurar el WiFi]
        Llama a \texttt{conectarWiFi()}. Si aun así no hay conexión, devuelve
        \texttt{false} de una vez: no tiene sentido intentar enviar sin red.
  \item[Líneas 5--7 --- preparar el cliente]
        \texttt{HTTPClient http;} crea el ``navegador''. Los dos \texttt{Timeout}
        ponen un límite de 2\,s para conectar y 2\,s para esperar respuesta. Si el
        PC está apagado o el firewall bloquea, la placa se rinde a los 2\,s en
        lugar de quedarse esperando mucho rato (y así puede seguir leyendo el
        botón).
  \item[Líneas 8--9 --- a dónde y en qué formato]
        \texttt{http.begin(PC\_URL)} fija la dirección de destino
        (\texttt{http://192.168.5.106:5000/alerta}). \texttt{addHeader} le avisa al
        servidor que los datos van en formato \textbf{JSON}.
  \item[Líneas 10--12 --- armar el mensaje]
        Se construye el texto JSON pegando pedazos con \texttt{+}. Las
        \texttt{\textbackslash"} son comillas que van \emph{dentro} del texto (sin la
        barra, C++ pensaría que el texto terminó). Pedazo a pedazo:
        \begin{center}\small
        \begin{tabular}{@{}ll@{}}
        \toprule
        \textbf{Pedazo del código} & \textbf{Queda en el mensaje}\\
        \midrule
        \texttt{"\{\textbackslash"dispositivo\textbackslash":\textbackslash"ESP32-S3\textbackslash",} & \texttt{\{"dispositivo":"ESP32-S3",}\\
        \texttt{\textbackslash"evento\textbackslash":\textbackslash""} & \texttt{"evento":"}\\
        \texttt{evento} & \texttt{ALERTA}\\
        \texttt{"\textbackslash",\textbackslash"ip\textbackslash":\textbackslash""} & \texttt{","ip":"}\\
        \texttt{WiFi.localIP().toString()} & \texttt{192.168.5.108}\\
        \texttt{"\textbackslash",\textbackslash"mac\textbackslash":\textbackslash""} & \texttt{","mac":"}\\
        \texttt{WiFi.macAddress()} & \texttt{E0:72:A1:AA:20:84}\\
        \texttt{"\textbackslash"\}"} & \texttt{"\}}\\
        \bottomrule
        \end{tabular}
        \end{center}
        Todo junto:
\begin{lstlisting}[numbers=none,language={}]
{"dispositivo":"ESP32-S3","evento":"ALERTA","ip":"192.168.5.108","mac":"E0:72:A1:AA:20:84"}
\end{lstlisting}
        La \textbf{MAC} es un identificador único de fábrica de la placa; sirve
        para distinguirla si algún día hay varias.
  \item[Línea 13 --- enviar]
        \texttt{http.POST(body)} envía el mensaje y \textbf{espera la respuesta} del
        servidor. El código de respuesta se guarda en \texttt{code}.
  \item[Línea 14 --- cerrar]
        \texttt{http.end()} cierra la conexión y libera memoria. Si no se cierra,
        la placa se va quedando sin memoria después de muchos envíos.
  \item[Línea 16 --- reportar]
        Imprime en el monitor algo como
        \texttt{Evento ALERTA enviado, respuesta HTTP: 200}. Aquí hay dos huecos:
        \texttt{\%s} (texto) se rellena con \texttt{evento} y \texttt{\%d} (número
        entero) con \texttt{code}.
  \item[Línea 17 --- ¿funcionó?]
        \texttt{code == 200} es \texttt{true} solo si el servidor contestó OK. Ese
        valor es el que se devuelve, y es el que usan \texttt{notificado} y
        \texttt{avisoConexion} para decidir si hay que reintentar.
  \item[Línea 20 --- atajo]
        \texttt{enviarNotificacion()} no hace nada nuevo: solo llama a
        \texttt{enviarEvento("ALERTA")}. Existe para que en el \texttt{loop()} se lea
        más claro qué se está enviando.
\end{description}

\subsection{\texttt{setup()}: lo que pasa al encender}

\begin{lstlisting}[language=C++]
void setup() {
  Serial.begin(115200);
  unsigned long t0 = millis();
  while (!Serial && millis() - t0 < 3000) delay(10);  // espera al monitor USB

  pinMode(PIN_BOTON, INPUT_PULLUP);
  pinMode(PIN_LED, OUTPUT);
  digitalWrite(PIN_LED, LOW);
  WiFi.mode(WIFI_STA);
  WiFi.setAutoReconnect(true);
  conectarWiFiConLed();
  Serial.println("Estado: ESTABLE");
}
\end{lstlisting}

\texttt{setup()} se ejecuta \textbf{una sola vez} al encender o al dar \textit{reset}:

\begin{enumerate}[leftmargin=*]
  \item Abre el puerto serial y espera hasta 3 s a que el monitor se conecte.
  \item Configura el botón como entrada con \textit{pull-up} y el LED como salida (apagado).
  \item \texttt{WIFI\_STA}: la ESP32 se conecta a una red existente (modo ``estación'',
        como un celular), no crea su propia red.
  \item \texttt{setAutoReconnect(true)}: si el WiFi se cae, se reconecta sola.
  \item Conecta al WiFi mostrando el estado con el LED.
\end{enumerate}

\subsection{\texttt{loop()}: el ciclo principal}

\texttt{loop()} se repite \textbf{miles de veces por segundo} mientras la placa esté
encendida. Cada vuelta es muy rápida porque nunca se usa \texttt{delay} largo.

\subsubsection*{Parte 1: leer el botón y avisar conexión}

\begin{lstlisting}[language=C++]
void loop() {
  bool presion = botonPresionado();

  // Aviso "CONECTADA" al PC: se reintenta cada 3 s hasta que el servidor responda
  if (!avisoConexion && estado == ESTABLE && WiFi.status() == WL_CONNECTED &&
      millis() - tAviso >= REINTENTO_MS) {
    tAviso = millis();
    avisoConexion = enviarEvento("CONECTADA");
  }
\end{lstlisting}

\begin{itemize}[leftmargin=*]
  \item Se lee el botón una vez por vuelta y se guarda en \texttt{presion}.
  \item Mientras no se haya confirmado el aviso de conexión, y la placa esté en
        reposo y con WiFi, se intenta enviar \texttt{"CONECTADA"} cada 3 s.
        Cuando el servidor responde OK, \texttt{avisoConexion} pasa a \texttt{true} y
        no se vuelve a enviar.
\end{itemize}

\subsubsection*{Parte 2: la máquina de estados}

\begin{lstlisting}[language=C++,firstnumber=10]
  switch (estado) {
    case ESTABLE:
      if (presion) {
        estado = PENDIENTE;
        tInicioAlerta = millis();
        digitalWrite(PIN_LED, HIGH);
        Serial.println("Estado: ALERTA (30 s para cancelar)");
      }
      break;
\end{lstlisting}

\textbf{ESTABLE:} si se oprime el botón, pasa a \texttt{PENDIENTE}, guarda la hora
exacta y enciende el LED.

\begin{lstlisting}[language=C++,firstnumber=19]
    case PENDIENTE:
      if (presion) {  // cancelado a tiempo
        estado = ESTABLE;
        digitalWrite(PIN_LED, LOW);
        Serial.println("Cancelado -> Estado: ESTABLE");
      } else if (millis() - tInicioAlerta >= VENTANA_MS) {
        estado = ENVIADA;
        notificado = false;
        primerIntento = true;
        Serial.println("30 s cumplidos -> enviando alerta");
      }
      break;
\end{lstlisting}

\textbf{PENDIENTE:} en cada vuelta se revisan dos cosas:
\begin{itemize}[leftmargin=*]
  \item Si se oprimió otra vez $\rightarrow$ fue un accidente: vuelve a
        \texttt{ESTABLE} y apaga el LED.
  \item Si ya pasaron 30 s desde \texttt{tInicioAlerta} $\rightarrow$ es una alerta
        real: pasa a \texttt{ENVIADA} y prepara el envío.
\end{itemize}

\begin{lstlisting}[language=C++,firstnumber=31]
    case ENVIADA:
      // LED queda encendido; se reintenta hasta que el PC confirme
      if (!notificado && (primerIntento || millis() - tUltimoIntento >= REINTENTO_MS)) {
        primerIntento = false;
        tUltimoIntento = millis();
        notificado = enviarNotificacion();
      }
      break;
  }
}
\end{lstlisting}

\textbf{ENVIADA:} el LED sigue encendido. Si la alerta aún no fue confirmada, se
envía: la primera vez de inmediato y, si falla, cada 3 s hasta que el servidor
responda OK. Después de eso el sistema se queda en este estado (el botón ya no hace
nada) hasta que se reinicie la placa.

%==========================================================
\section{\texttt{servidor\_pc.py}: el receptor en el computador}
%==========================================================

Es un servidor web mínimo hecho solo con la librería estándar de Python (no hay que
instalar nada).

\subsection{Importaciones y puerto}

\begin{lstlisting}[language=Python]
import json
import socket
from datetime import datetime
from http.server import BaseHTTPRequestHandler, HTTPServer

PUERTO = 5000
\end{lstlisting}

\begin{itemize}[leftmargin=*]
  \item \texttt{json}: convertir el texto recibido en un diccionario de Python.
  \item \texttt{socket}: averiguar la IP del computador.
  \item \texttt{datetime}: poner fecha y hora a cada mensaje.
  \item \texttt{http.server}: el servidor web.
  \item \texttt{PUERTO = 5000}: la ``puerta'' por donde entra la conexión. Debe
        coincidir con el \texttt{:5000} de \texttt{PC\_URL}.
\end{itemize}

\subsection{El manejador de peticiones}

\begin{idea}[¿Qué es una clase y qué es \texttt{self}?]
Una \textbf{clase} es como un molde que agrupa funciones relacionadas. Aquí la
clase \texttt{Handler} (``manejador'') reúne todo lo que el servidor hace cuando
llega un mensaje. \texttt{(BaseHTTPRequestHandler)} significa que la clase
\textbf{hereda} de un manejador que ya trae Python: ya sabe leer mensajes HTTP y
responder, y nosotros solo escribimos qué hacer con ellos.

\texttt{self} es ``este mensaje que acaba de llegar''. A través de \texttt{self}
se accede a todo lo del mensaje: la ruta (\texttt{self.path}), los encabezados
(\texttt{self.headers}), quién lo envió (\texttt{self.client\_address}), etc.
\end{idea}

\subsubsection*{Parte 1: recibir y leer el mensaje}

\begin{lstlisting}[language=Python,firstnumber=13]
class Handler(BaseHTTPRequestHandler):
    def do_POST(self):
        if self.path != "/alerta":
            self.send_response(404)
            self.end_headers()
            return
        largo = int(self.headers.get("Content-Length", 0))
        datos = self.rfile.read(largo).decode(errors="ignore")
        try:
            datos = json.loads(datos)
        except json.JSONDecodeError:
            pass
\end{lstlisting}

\begin{description}[leftmargin=*,style=nextline]
  \item[Línea 14 --- \texttt{do\_POST}]
        El nombre no es casual: Python llama \textbf{automáticamente} a la función
        \texttt{do\_POST} cada vez que llega un POST. Es la pareja del
        \texttt{http.POST(body)} de la ESP32. Nadie la llama a mano.
  \item[Líneas 15--18 --- revisar la dirección]
        \texttt{self.path} es la ruta a la que llegó el mensaje. Si no es
        \texttt{/alerta} (por ejemplo, alguien escribió mal la URL), responde
        \texttt{404} (``no existe''), cierra los encabezados de la respuesta y
        \texttt{return} termina sin hacer nada más.
  \item[Línea 19 --- ¿qué tan largo es el mensaje?]
        Todo mensaje HTTP trae un encabezado \texttt{Content-Length} que dice
        cuántos bytes mide el contenido. \texttt{.get("Content-Length", 0)} lo lee
        (y si no viene, usa 0). \texttt{int(...)} lo convierte de texto a número.
  \item[Línea 20 --- leer el contenido]
        \texttt{self.rfile} es el ``archivo'' por donde llegan los datos.
        \texttt{.read(largo)} lee exactamente esa cantidad de bytes, y
        \texttt{.decode(...)} los convierte de bytes a texto normal.
        \texttt{errors="ignore"} evita que el servidor se caiga si viene algún
        carácter raro. Ahora \texttt{datos} es un texto como
        \texttt{\{"evento":"ALERTA", ...\}}.
  \item[Líneas 21--24 --- convertir el JSON]
        \texttt{json.loads} convierte ese texto en un \textbf{diccionario} de
        Python, con el que se puede preguntar \texttt{datos["evento"]} y obtener
        \texttt{"ALERTA"}. El \texttt{try/except} es una red de seguridad: si el
        texto no es JSON válido, en vez de caerse el programa, sigue
        (\texttt{pass} = no hacer nada) y \texttt{datos} se queda como texto.
\end{description}

\subsubsection*{Parte 2: decidir qué mostrar y responder}

\begin{lstlisting}[language=Python,firstnumber=25]
        hora = datetime.now().strftime("%Y-%m-%d %H:%M:%S")
        evento = datos.get("evento") if isinstance(datos, dict) else None
        if evento == "CONECTADA":
            print(f"[{hora}] ESP32 conectada al WiFi desde {self.client_address[0]}")
        else:
            print("\a")  # pitido de consola
            print("=" * 50)
            print(f"  ¡ALERTA RECIBIDA!  {hora}")
            print(f"  Desde: {self.client_address[0]}  Datos: {datos}")
            print("=" * 50)
        self.send_response(200)
        self.end_headers()
        self.wfile.write(b"OK")
\end{lstlisting}

\begin{description}[leftmargin=*,style=nextline]
  \item[Línea 25 --- la hora]
        \texttt{datetime.now()} da la fecha y hora actuales, y \texttt{strftime}
        les da formato. Cada código se reemplaza por un dato: \texttt{\%Y} año,
        \texttt{\%m} mes, \texttt{\%d} día, \texttt{\%H} hora, \texttt{\%M}
        minutos, \texttt{\%S} segundos. Resultado: \texttt{2026-09-25 21:32:19}.
  \item[Línea 26 --- ¿qué evento es?]
        También es un ternario, en su versión de Python:
        \begin{center}
        \texttt{valor\_si\_verdad if condición else valor\_si\_falso}
        \end{center}
        \texttt{isinstance(datos, dict)} pregunta si \texttt{datos} sí quedó
        convertido en diccionario. Si sí, saca el campo \texttt{"evento"}; si no
        (el JSON venía dañado), usa \texttt{None} (``nada'').
  \item[Líneas 27--28 --- mensaje de conexión]
        Si el evento es \texttt{"CONECTADA"}, imprime una sola línea. La
        \texttt{f} antes de las comillas crea un \textbf{f-string}: lo que va entre
        llaves \texttt{\{ \}} se reemplaza por su valor.
        \texttt{self.client\_address[0]} es la IP de quien envió el mensaje, o sea,
        la placa. Resultado:
        \texttt{[2026-09-25 21:32:19] ESP32 conectada al WiFi desde 192.168.5.108}.
  \item[Líneas 29--34 --- mensaje de alerta]
        Cualquier otro evento se trata como alerta y se muestra llamativo:
        \begin{itemize}
          \item \texttt{print("\textbackslash a")}: \texttt{\textbackslash a} es un carácter especial
                que hace sonar un pitido en la terminal.
          \item \texttt{"=" * 50}: en Python, multiplicar un texto lo repite; esto
                dibuja una línea de 50 signos \texttt{=}.
          \item Luego imprime el título con la hora, y la IP y los datos que llegaron.
        \end{itemize}
  \item[Líneas 35--37 --- responder OK]
        \texttt{send\_response(200)} arma la respuesta ``recibido OK'',
        \texttt{end\_headers()} cierra los encabezados, y
        \texttt{wfile.write(b"OK")} escribe el contenido (la \texttt{b} indica que
        va en bytes). \textbf{Esto es clave}: ese \texttt{200} es exactamente el
        \texttt{code} que recibe la ESP32 en la línea \texttt{int code = http.POST(body)}.
        Si no llega, la placa asume que falló y reintenta cada 3\,s.
\end{description}

\subsubsection*{Parte 3: silenciar el registro automático}

\begin{lstlisting}[language=Python,firstnumber=39]
    def log_message(self, *args):
        pass  # silencia el log por defecto
\end{lstlisting}

Por defecto, \texttt{BaseHTTPRequestHandler} imprime una línea técnica por cada
mensaje que llega, algo como:
\begin{lstlisting}[numbers=none,language={}]
192.168.5.108 - - [25/Sep/2026 21:32:19] "POST /alerta HTTP/1.1" 200 -
\end{lstlisting}
Para hacerlo, llama internamente a una función llamada \texttt{log\_message}. Al
escribir nuestra propia \texttt{log\_message} dentro de \texttt{Handler}, la
\textbf{reemplazamos}. La nuestra solo tiene \texttt{pass} (no hace nada), así que
esas líneas desaparecen y en la terminal solo quedan nuestros mensajes, limpios.
\texttt{*args} significa ``acepta cualquier cantidad de datos'', porque Python le
pasa varios y no nos interesa ninguno.

\subsection{Arranque del servidor}

\begin{lstlisting}[language=Python,firstnumber=43]
if __name__ == "__main__":
    s = socket.socket(socket.AF_INET, socket.SOCK_DGRAM)
    try:
        s.connect(("8.8.8.8", 80))
        ip = s.getsockname()[0]
    except OSError:
        ip = "desconocida"
    finally:
        s.close()
    print(f"Escuchando en http://{ip}:{PUERTO}/alerta")
    print("Pon esa IP en PC_URL del sketch del ESP32.")
    HTTPServer(("0.0.0.0", PUERTO), Handler).serve_forever()
\end{lstlisting}

\begin{description}[leftmargin=*,style=nextline]
  \item[Línea 43 --- punto de arranque]
        \texttt{if \_\_name\_\_ == "\_\_main\_\_":} significa ``ejecuta esto solo si
        el archivo se corrió directamente'' (con \texttt{python pc/servidor\_pc.py}).
        Es la forma estándar de marcar dónde empieza un programa en Python.
  \item[Líneas 44--51 --- averiguar la IP del PC]
        Un computador puede tener varias IP (WiFi, cable, VPN\ldots). Para saber
        cuál usa en la red local se hace un truco:
        \begin{enumerate}
          \item Se crea una ``conexión'' de red tipo UDP con
                \texttt{socket.socket(AF\_INET, SOCK\_DGRAM)} (\texttt{AF\_INET} = direcciones IPv4).
          \item \texttt{s.connect(("8.8.8.8", 80))} la apunta hacia un servidor de
                Google. Con UDP esto \textbf{no envía nada}; solo obliga al sistema a
                decidir por cuál tarjeta de red saldría.
          \item \texttt{s.getsockname()[0]} pregunta qué IP local eligió. Esa es la
                IP del PC en tu red (\texttt{192.168.5.106}).
          \item \texttt{except OSError}: si no hay red, en vez de caerse pone
                \texttt{"desconocida"}.
          \item \texttt{finally}: pase lo que pase, \texttt{s.close()} cierra la
                conexión.
        \end{enumerate}
  \item[Líneas 52--53 --- mostrar la dirección]
        Imprime la URL completa, que es justo lo que hay que copiar en
        \texttt{PC\_URL} de \texttt{secrets.h}.
  \item[Línea 54 --- arrancar y quedarse escuchando]
        Esta línea hace tres cosas:
        \begin{itemize}
          \item \texttt{("0.0.0.0", PUERTO)}: escuchar en el puerto 5000 por
                \textbf{todas} las tarjetas de red del PC, para que la ESP32 lo
                encuentre sin importar por cuál llegue.
          \item \texttt{Handler}: cada mensaje que llegue se atiende con nuestra
                clase (la que tiene \texttt{do\_POST}).
          \item \texttt{.serve\_forever()}: se queda esperando mensajes
                \textbf{para siempre}. Por eso la terminal parece ``congelada'' con
                el cursor: está escuchando. Se detiene con \texttt{Ctrl+C}.
        \end{itemize}
\end{description}

%==========================================================
\section{Flujo completo}
%==========================================================

\begin{center}
\begin{tikzpicture}[
  font=\small, >={Stealth[length=2.5mm]},
  msg/.style={->,thick}, ev/.style={align=left,font=\footnotesize}]
  \node (U) at (0,0) {\textbf{Usuario}};
  \node (E) at (4.5,0) {\textbf{ESP32-S3}};
  \node (P) at (9,0) {\textbf{PC (servidor)}};
  \foreach \x in {0,4.5,9} \draw[gray,dashed] (\x,-0.3) -- (\x,-9.4);
  \draw[msg] (4.5,-0.8) -- node[above]{POST \texttt{CONECTADA}} (9,-0.8);
  \draw[msg] (9,-1.3) -- node[above]{\texttt{200 OK}} (4.5,-1.3);
  \node[ev,right] at (9.1,-1.05) {``ESP32 conectada''};
  \draw[msg] (0,-2.2) -- node[above]{oprime botón} (4.5,-2.2);
  \node[ev,left] at (4.4,-2.7) {LED on, cuenta 30 s};
  \draw[msg] (0,-3.4) -- node[above]{oprime otra vez} (4.5,-3.4);
  \node[ev,left] at (4.4,-3.9) {cancelado, LED off};
  \draw[gray] (-0.4,-4.4) -- (9.6,-4.4);
  \draw[msg] (0,-5.2) -- node[above]{oprime botón} (4.5,-5.2);
  \node[ev,left] at (4.4,-5.7) {LED on, cuenta 30 s};
  \draw[thick,gray,dotted] (4.7,-5.9) -- (4.7,-7.2);
  \node[ev,right] at (4.8,-6.55) {pasan 30 s};
  \draw[msg] (4.5,-7.6) -- node[above]{POST \texttt{ALERTA}} (9,-7.6);
  \node[ev,right] at (9.1,-7.85) {``¡ALERTA RECIBIDA!''};
  \draw[msg] (9,-8.1) -- node[above]{\texttt{200 OK}} (4.5,-8.1);
  \node[ev,left] at (4.4,-8.8) {LED queda on};
\end{tikzpicture}
\end{center}

%==========================================================
\section{Uso y solución de problemas}
%==========================================================

\subsection{Pasos de uso}
\begin{enumerate}[leftmargin=*]
  \item En el PC: \texttt{python pc/servidor\_pc.py} y dejar esa terminal abierta.
  \item Poner la IP que muestra el servidor en \texttt{PC\_URL} de \texttt{secrets.h}.
  \item Subir el firmware con \textbf{Upload} ($\rightarrow$) en PlatformIO.
  \item Dar \textit{reset} a la placa y comprobar que aparece ``ESP32 conectada''.
\end{enumerate}

\subsection{Problemas comunes}

\begin{center}
\begin{tabular}{@{}p{5.2cm}p{9.5cm}@{}}
\toprule
\textbf{Síntoma} & \textbf{Qué revisar}\\
\midrule
LED parpadea rápido al arrancar & Nombre o clave del WiFi en \texttt{secrets.h}; que la red sea de 2.4\,GHz.\\
LED fijo 2 s pero no llega nada al PC & Que el servidor esté corriendo; firewall de Windows (puerto 5000); que la IP del PC no haya cambiado.\\
\texttt{Could not open COM..} al subir & Cable USB conectado; escoger el puerto correcto; modo carga manual: mantener BOOT, pulsar RST, soltar BOOT.\\
Monitor serial en blanco & Normal si se abrió después de arrancar; dar RST con el monitor abierto.\\
Cambiar de computador & Correr el servidor en el otro PC, poner su IP en \texttt{PC\_URL} y volver a subir el código.\\
\bottomrule
\end{tabular}
\end{center}

\end{document}
